% ============================================================================
% หัวข้อใหม่สำหรับบทที่ 4 : การทดลองที่ 1 (ฉบับย่อ)
% ต้องมีใน preamble : \usepackage{booktabs,graphicx,amsmath,amssymb,array}
% และ environment proposition (ข้อเสนอ) แบบเดียวกับข้อเสนอ 2.1
% ตารางและรูปสร้างอัตโนมัติจาก  python run_exp1.py   ไว้ที่ results/real/
% ข้อความใน [วงเล็บเหลี่ยม] ต้องเขียนหลังเห็นผลจริงเท่านั้น
% ============================================================================

\section{การทดลองที่ 1: เปรียบเทียบ Ambiguity Set ภายใต้ตัวพยากรณ์เดียวกัน}
\label{sec:exp1}

การทดลองนี้ตอบคำถามว่า เมื่อใช้ตัวพยากรณ์ AR(1)--GARCH(1,1)-t ชุดเดียวกัน การเลือก
Ambiguity Set แบบ Nominal, Wasserstein และ Possibility ให้ผลต่อแนวหน้าพาเรโตและผลลัพธ์
นอกกลุ่มตัวอย่างต่างกันอย่างไร โดยควบคุมให้ทุกแบบจำลองใช้ค่าพยากรณ์ สถานการณ์จำลอง
ตัวแก้ปัญหา ข้อจำกัด และวันปรับพอร์ตเดียวกันทั้งหมด ความแตกต่างของผลจึงมาจาก Ambiguity
Set เพียงอย่างเดียว ส่วนการเรียนรู้ร่วม (Decision-aware Learning) จะทดสอบในการทดลองที่ 2

\subsection{การออกแบบเพื่อลดความเอนเอียงของผลการทดลอง}
\label{sec:exp1-bias}

เนื่องจากช่องว่างงานวิจัยข้อที่ 4 คือการประเมินผลที่ขาดความเข้มงวด การทดลองนี้จึงใช้เกณฑ์เดียวกัน
กับที่ใช้ประเมินแบบจำลอง E-DRO ในหัวข้อที่ 4.2--4.6 กับงานของโครงการเอง ดังตารางที่~\ref{tab:exp1-bias}
การตั้งค่าทุกรายการกำหนดไว้ในไฟล์ \texttt{config.py} ก่อนรันกับข้อมูลทดสอบ และรายงานรหัส
Commit พร้อมผลการทดลอง

\begin{table}[htbp]
\centering
\caption{แหล่งของความเอนเอียงและวิธีป้องกันในการทดลองที่ 1}
\label{tab:exp1-bias}
\begin{tabular}{>{\raggedright\arraybackslash}p{0.32\linewidth}>{\raggedright\arraybackslash}p{0.6\linewidth}}
\toprule
แหล่งของความเอนเอียง & วิธีป้องกัน \\
\midrule
การเลือกสินทรัพย์ที่ทราบภายหลังว่าให้ผลดี (Survivorship / Hindsight Bias)
 & ใช้กองทุนดัชนีรายกลุ่มอุตสาหกรรม 9 กองทุน หนึ่งกองทุนต่อหนึ่ง GICS Sector ซึ่งซื้อขายต่อเนื่องตลอดช่วงข้อมูล \\
การเลือกช่วงเวลาทดสอบ
 & ประเมินนอกกลุ่มตัวอย่างต่อเนื่องตลอดช่วงข้อมูลที่มี และแยกผลตามช่วงตลาดที่กำหนดล่วงหน้า \\
การใช้ข้อมูลอนาคต (Look-ahead Bias)
 & ทุกการตัดสินใจ ณ สัปดาห์ $t$ ใช้ข้อมูลถึงสัปดาห์ $t$ เท่านั้น และมีการทดสอบในโปรแกรม \\
การปรับพารามิเตอร์จนได้ผลดี (Data Snooping)
 & $\varepsilon$, $\kappa$, $\lambda$ และ $\alpha$ กำหนดด้วยกฎล่วงหน้า และรายงานทุกค่าที่ทดลอง \\
การเปรียบเทียบที่ไม่เท่าเทียม
 & ใช้ตัวพยากรณ์ สถานการณ์จำลอง และตัวแก้ปัญหาเดียวกัน และเปรียบเทียบกับพอร์ต $1/N$ ทุกครั้ง \\
ความแตกต่างที่เกิดจากความบังเอิญ
 & ทดสอบ Sharpe Ratio ด้วยวิธีของ Ledoit และ Wolf ปรับค่า $p$ ด้วยวิธี Holm และใช้ Block Bootstrap กับ Hypervolume \\
ต้นทุนการซื้อขาย & รายงานผลทั้งก่อนและหลังหักต้นทุน 10 Basis Points ต่อการซื้อขาย \\
\bottomrule
\end{tabular}
\end{table}

\subsection{ข้อมูลและการตั้งค่า}

ใช้ราคาปิดที่ปรับปรุงแล้วของกองทุน Select Sector SPDR จำนวน 9 กองทุน ได้แก่ XLB (Materials),
XLE (Energy), XLF (Financials), XLI (Industrials), XLK (Information Technology), XLP (Consumer
Staples), XLU (Utilities), XLV (Health Care) และ XLY (Consumer Discretionary) ซึ่งจัดกลุ่มตาม
GICS ไม่รวม XLRE และ XLC เพราะเริ่มซื้อขายในปี ค.ศ. 2015 และ 2018 ตามลำดับ ทำให้ข้อมูลสั้น
เกินไป ข้อมูลรายสัปดาห์ใช้ราคาวันทำการสุดท้ายของแต่ละสัปดาห์ และคำนวณผลตอบแทนแบบลอการิทึม
ตาม (3.1)

ควรระบุไว้ด้วยว่า โครงสร้าง GICS มีการปรับปรุงระหว่างช่วงข้อมูล ได้แก่ การแยก Real Estate ออก
จาก Financials ในปี ค.ศ. 2016 การจัดตั้ง Communication Services ในปี ค.ศ. 2018 ซึ่งย้ายบริษัท
อินเทอร์เน็ตบางส่วนออกจาก Information Technology และการปรับโครงสร้างในเดือนมีนาคม ค.ศ. 2023
ส่วนประกอบของกองทุนจึงเปลี่ยนไปตามเวลา ซึ่งเป็นการเปลี่ยนแปลงของการแจกแจง (Distributional
Shift) ที่เกิดขึ้นจริงในตลาดและสอดคล้องกับเหตุผลของการใช้ DRO

ในแต่ละสัปดาห์ $t$ ใช้หน้าต่างข้อมูลย้อนหลัง $260$ สัปดาห์ ประมาณพารามิเตอร์ของ AR(1)--GARCH(1,1)-t
ใหม่ทุก $4$ สัปดาห์ และกรองความผันผวนด้วยพารามิเตอร์ล่าสุดทุกสัปดาห์ จากนั้นเลือกน้ำหนักพอร์ต
$w_t$ และถือครองตลอดสัปดาห์ $t+1$ ผลนอกกลุ่มตัวอย่างจึงครอบคลุมตั้งแต่สัปดาห์ที่ $261$ ของข้อมูล
จนถึงสัปดาห์สุดท้าย การตั้งค่าทั้งหมดแสดงในตารางที่~\ref{tab:exp1-config}

\begin{table}[htbp]
\centering
\caption{การตั้งค่าของการทดลองที่ 1 (กำหนดก่อนการรันกับข้อมูลทดสอบ)}
\label{tab:exp1-config}
\begin{tabular}{ll}
\toprule
รายการ & ค่าที่ใช้ \\
\midrule
สินทรัพย์ & Sector ETF 9 กองทุน (GICS) \\
หน้าต่างข้อมูลฝึก / การประมาณใหม่ & 260 สัปดาห์ / ทุก 4 สัปดาห์ \\
ความถี่การปรับพอร์ต & ทุกสัปดาห์ \\
จำนวนสถานการณ์จำลอง $N$ & 259 (FHS จาก Standardized Residuals) \\
ระดับ CVaR & $\alpha = 0.05$ \\
น้ำหนักวัตถุประสงค์ & $\lambda \in \{0, 0.1, \ldots, 1\}$; ตารางผลตอบแทนใช้ $\lambda = 0.5$ \\
รัศมี Wasserstein & $\varepsilon = c\,\varepsilon_0$, $c \in \{0.2, 1\}$ (ภาคผนวก: $0.5, 2$) \\
ระดับของ Possibility (วิธี B) & $\gamma \in \{0.05, 0.10, \ldots, 0.95\}$ \\
พจน์ Ridge & $10^{-6}\lVert w \rVert_2^2$ (ให้คำตอบมีเพียงคำตอบเดียว) \\
ต้นทุนการซื้อขาย & 0 และ 10 Basis Points \\
\bottomrule
\end{tabular}
\end{table}

\subsection{สถานการณ์จำลองร่วมจากตัวพยากรณ์}

ให้ $\hat\mu_{t+1}, \hat\sigma_{t+1} \in \mathbb{R}^n$ เป็นค่าเฉลี่ยและส่วนเบี่ยงเบนมาตรฐานแบบมีเงื่อนไขที่พยากรณ์
หนึ่งขั้น และ $\hat z_s \in \mathbb{R}^n$, $s = 1,\ldots,N$ เป็นเวกเตอร์ Standardized Residuals ของ
สัปดาห์เดียวกันในหน้าต่างข้อมูลฝึก สถานการณ์จำลองแบบ Filtered Historical Simulation คือ
\begin{equation}
\xi_s = \hat\mu_{t+1} + \hat\sigma_{t+1} \odot \hat z_s, \qquad
L_s(w) = -\xi_s^\top w, \qquad s = 1, \ldots, N
\label{eq:exp1-fhs}
\end{equation}
การใช้ Residuals ทั้งเวกเตอร์ของสัปดาห์เดียวกันรักษาสหสัมพันธ์ระหว่างสินทรัพย์และหางหนาเชิง
ประจักษ์ไว้ โดยไม่ต้องสมมติรูปแบบของการแจกแจงร่วม ทุกแบบจำลองในการทดลองนี้ใช้
สถานการณ์จำลองชุดเดียวกัน ยกเว้น Possibility วิธี A ซึ่งใช้การแจกแจงแบบวงรีตามบทแทรก 2.1
เพื่อเป็นกรณีเปรียบเทียบ

\subsection{แบบจำลองที่เปรียบเทียบ}

ทุกแบบจำลองแก้ปัญหา
\begin{equation}
w^\ast_\lambda \in \arg\min_{w \in \mathcal{X}} \; \lambda F_1(w) + (1-\lambda) F_2(w) + 10^{-6}\lVert w \rVert_2^2
\label{eq:exp1-ws}
\end{equation}
โดย $F_1$ คือค่าคาดหวังของผลขาดทุนในกรณีเลวร้ายที่สุด และ $F_2$ คือ CVaR$_\alpha$ ในกรณีเลวร้ายที่สุด
เนื่องจากทุกปัญหาเป็นปัญหานูนตามข้อเสนอ 2.1 การกวาดค่า $\lambda$ จึงได้แนวหน้าพาเรโตครบถ้วน

\paragraph{Nominal}
$F_1 = \frac{1}{N}\sum_s L_s(w)$ และ $F_2$ คือ CVaR เชิงประจักษ์ในรูป LP (2.16)

\paragraph{Wasserstein}
กำหนด Ambiguity Set เป็นลูกบอล $W_1$ รัศมี $\varepsilon$ รอบการแจกแจงเชิงประจักษ์ของ Residuals
ในปริภูมิ Residual ด้วย Ground Metric แบบ $\ell_2$ และไม่จำกัด Support เนื่องจาก
$L(w, z) = -\hat\mu^\top w - (\hat\sigma \odot w)^\top z$ เป็นฟังก์ชันอฟฟีนใน $z$ ที่มีค่าคงที่ลิปชิตซ์
$\lVert \hat\sigma \odot w \rVert_2$ ผลของ Mohajerin Esfahani และ Kuhn [33] ให้
\begin{equation}
F_1^{W} = F_1^{\mathrm{nom}} + \varepsilon \lVert \hat\sigma \odot w \rVert_2, \qquad
F_2^{W} = F_2^{\mathrm{nom}} + \frac{\varepsilon}{\alpha} \lVert \hat\sigma \odot w \rVert_2
\label{eq:exp1-wass}
\end{equation}
การกำหนดลูกบอลในปริภูมิ Residual ทำให้ $\varepsilon$ ไม่มีหน่วย และพจน์ความทนทานปรับตาม
ความผันผวนที่พยากรณ์ได้ในแต่ละสัปดาห์ การเลือก $\ell_2$ แทน $\ell_\infty$ หลีกเลี่ยงกรณีที่นอร์มคู่
เป็นค่าคงที่บน Simplex ตามที่ Yadav และ Mehra [39] ชี้ไว้ (หมายเหตุ 2.4)

รัศมีฐาน $\varepsilon_0$ คำนวณใหม่ทุกครั้งที่ประมาณพารามิเตอร์ จากข้อมูลฝึกเท่านั้น โดยสร้าง
Stationary Bootstrap ของเวกเตอร์ Residuals (ความยาวบล็อกเฉลี่ย 4 สัปดาห์) จำนวน 100 ชุด
คำนวณ $W_1$ แบบแม่นตรงระหว่างแต่ละชุดกับข้อมูลเดิมด้วยปัญหาการจับคู่ (2.36) และใช้ควอนไทล์ที่
$95\%$ เป็น $\varepsilon_0$ จากนั้นรายงานผลที่ $\varepsilon = c\,\varepsilon_0$ ทุกค่า $c$ ที่กำหนดไว้

\paragraph{Possibility วิธี A (การแจกแจงแบบวงรี)}
สมมติให้ Residual มีการแจกแจง Student's $t$ หลายตัวแปรแบบปรับมาตรฐาน โดยมีเมทริกซ์สหสัมพันธ์
$S$ และองศาอิสระ $\nu$ (ค่ามัธยฐานของ $\nu_i$ ที่ประมาณได้) การแปลงความน่าจะเป็นเป็นความเป็นไป
ได้ให้ $\alpha$-cut เป็นทรงรีซ้อนกัน และ Ambiguity Set คือการแจกแจงที่รัศมีมหาลาโนบิส
$\rho = \sqrt{z^\top S^{-1} z}$ มีค่าน้อยกว่าหรือเท่ากับรัศมีอ้างอิง $R$ ในเชิงสโตแคสติก เนื่องจาก
$\max\{-z^\top(\hat\sigma\odot w) : \sqrt{z^\top S^{-1}z} = \rho\} = \rho\,\lVert S^{1/2}(\hat\sigma \odot w)\rVert_2$
การแจกแจงที่เลวร้ายที่สุดจึงวางมวลทั้งหมดในทิศทางเดียวกันตามรัศมี และได้
\begin{equation}
F_1^{A} = -\hat\mu^\top w + \kappa_E \lVert S^{1/2} (\hat\sigma \odot w) \rVert_2, \qquad
F_2^{A} = -\hat\mu^\top w + \kappa_C \lVert S^{1/2} (\hat\sigma \odot w) \rVert_2
\label{eq:exp1-possA}
\end{equation}
โดย $\kappa_E = \mathbb{E}[R]$ และ $\kappa_C = \mathrm{CVaR}_\alpha(R)$ ซึ่งคำนวณจาก
$R^2 = \tfrac{\nu-2}{\nu} n \, F(n, \nu)$ ทั้งสองวัตถุประสงค์มีรูปเดียวกันและต่างกันเพียงค่าคงที่ ดังนั้น
(\ref{eq:exp1-ws}) จึงเป็นปัญหา Mean--Standard Deviation ที่มีสัมประสิทธิ์
$\lambda\kappa_E + (1-\lambda)\kappa_C$ แนวหน้าพาเรโตจึงเป็นเพียงช่วงสั้น ๆ ของ Mean--Standard
Deviation Frontier ตามที่ระบุในหัวข้อที่ 5.2 วิธี A จึงใช้เพื่อแสดงปรากฏการณ์นี้เชิงประจักษ์

\paragraph{Possibility วิธี B (สถานการณ์จำลอง)}
ใช้สถานการณ์จำลอง (\ref{eq:exp1-fhs}) ชุดเดียวกับแบบจำลองอื่น เรียงลำดับ $\hat z_s$ ตามระยะทาง
มหาลาโนบิสจากจุดพยากรณ์ และกำหนดความเป็นไปได้เชิงประจักษ์ $\pi_s$ เป็นสัดส่วนของสถานการณ์
ที่อยู่ไกลเท่ากันหรือไกลกว่า สำหรับแต่ละระดับ $\gamma_k$ ให้ $C_k$ เป็นเซตของสถานการณ์ที่อยู่ใกล้ที่สุด
$\lceil (1-\gamma_k) N \rceil$ จุด การแจกแจง $p$ บนสถานการณ์จำลองอยู่ใน Ambiguity Set เมื่อ
$p(C_k) \ge 1 - \gamma_k$ ทุก $k$ ซึ่งเป็นรูปเชิงประจักษ์ของเงื่อนไข $P(A) \le \Pi(A)$ เขียนเป็น
$A p \ge b$ โดย $A_{ks} = \mathbb{1}[s \in C_k]$ และ $b_k = 1-\gamma_k$

\begin{proposition}
\label{prop:exp1-possB}
สำหรับ $w$ คงที่ ค่าในกรณีเลวร้ายที่สุดของ Possibility วิธี B เท่ากับค่าเหมาะที่สุดของ LP ต่อไปนี้
\begin{align}
F_1^{B}(w) &= \min_{t,\; y \ge 0} \; t - b^\top y
  \quad \text{s.t.} \quad t - (A^\top y)_s \ge L_s(w) \;\; \forall s \label{eq:exp1-possB1}\\
F_2^{B}(w) &= \min_{\eta,\; t,\; y \ge 0} \; \eta + \tfrac{1}{\alpha}(t - b^\top y)
  \quad \text{s.t.} \quad t - (A^\top y)_s \ge \max\{L_s(w) - \eta, 0\} \;\; \forall s \label{eq:exp1-possB2}
\end{align}
\end{proposition}

\noindent\textit{แนวการพิสูจน์.} $F_1^B = \max\{L^\top p : \mathbf{1}^\top p = 1,\, Ap \ge b,\, p \ge 0\}$ เป็น LP
ที่มีคำตอบเป็นไปได้ (การแจกแจงเชิงประจักษ์ $p = \mathbf{1}/N$ สอดคล้องกับทุกเงื่อนไข) และมีขอบเขต
ภาวะคู่กันอย่างเข้มของ LP ให้ (\ref{eq:exp1-possB1}) สำหรับ $F_2^B$ ใช้สูตร (2.15) แล้วสลับลำดับ
$\max_p$ และ $\min_\eta$ ได้ เพราะฟังก์ชันนูนใน $\eta$ และเชิงเส้นใน $p$ บนเซตที่เป็น Compact
จากนั้นใช้ภาวะคู่กันกับปัญหาชั้นในตามเดิม \hfill$\square$

เนื่องจาก $L_s(w)$ เป็นเชิงเส้นใน $w$ ปัญหา (\ref{eq:exp1-ws}) ของวิธี B จึงเป็น LP เดียวในตัวแปร
$(w, t, y, \eta)$ และค่าเฉลี่ยกับส่วนหางถูกทำให้เลวร้ายด้วยกลไกต่างกัน วัตถุประสงค์ทั้งสองจึงไม่
ซ้อนกันเหมือนวิธี A ข้อจำกัดของวิธีนี้คือ Support ถูกจำกัดอยู่ที่สถานการณ์จำลอง ค่าที่ได้จึงเป็นค่า
ประมาณเชิงข้อมูลของ Possibility-induced Ambiguity Set และอาจต่ำกว่าค่าในกรณีต่อเนื่อง

\subsection{การประเมินผล}

สำหรับแต่ละแบบจำลองและแต่ละ $\lambda$ คำนวณผลตอบแทนจริงของพอร์ตด้วย Simple Return ตาม (3.22)
และสร้างแนวหน้านอกกลุ่มตัวอย่าง โดยให้แต่ละ $\lambda$ เป็นหนึ่งจุด
$(\overline{\ell}, \widehat{\mathrm{CVaR}}_\alpha)$ ของผลขาดทุนที่เกิดขึ้นจริงตลอดช่วงทดสอบ
Hypervolume คำนวณหลังปรับมาตราส่วนแบบ Min--Max ด้วยค่าจากทุกแบบจำลองร่วมกัน และใช้จุด
อ้างอิง $(1.1, 1.1)$ เดียวกัน ช่วงความเชื่อมั่นของผลต่าง HV เทียบกับ Nominal ใช้ Stationary Block
Bootstrap ของสัปดาห์ทดสอบ (ความยาวบล็อกเฉลี่ย 8 สัปดาห์, 500 ชุด) โดยสุ่มสัปดาห์ชุดเดียวกันให้ทุก
แบบจำลอง Sharpe Ratio ทดสอบด้วยวิธี HAC ของ Ledoit และ Wolf [56] และปรับค่า $p$ ด้วยวิธี
Holm ภายในกลุ่มการเปรียบเทียบเดียวกัน

ข้อควรระวังในการตีความ HV คือ ค่านี้ให้รางวัลแก่แนวหน้าที่ยาว แบบจำลองที่ทนทานมากจะเลือก
พอร์ตที่ใกล้กันในทุกค่า $\lambda$ จึงได้แนวหน้าสั้นและอาจได้ HV ต่ำ แม้จุดของแบบจำลองนั้นจะอยู่ต่ำกว่า
แนวหน้าของ Nominal ในช่วงที่ซ้อนทับกัน การสรุปผลจึงต้องพิจารณา HV ร่วมกับรูปที่~\ref{fig:exp1-front}
และตัวชี้วัดราย $\lambda$ ในตารางที่~\ref{tab:exp1-perf} เสมอ ช่วงตลาดผันผวนกำหนดล่วงหน้า ได้แก่ GFC (10/2007--03/2009),
Euro Debt (07--12/2011), COVID (02--06/2020) และ Bear 2022 (01--10/2022)

\subsection{การตรวจสอบความถูกต้องของโปรแกรม}

ก่อนรันกับข้อมูลจริง โปรแกรมผ่านการตรวจสอบทั้งหมดใน \texttt{tests.py} ได้แก่
(1) ค่าพยากรณ์ไม่เปลี่ยนเมื่อแก้ไขข้อมูลหลังสัปดาห์ที่ตัดสินใจ (ไม่มี Look-ahead)
(2) Wasserstein ที่ $\varepsilon = 0$ ให้ค่าเท่ากับ Nominal และค่าในกรณีเลวร้ายที่สุดไม่ลดลงเมื่อ $\varepsilon$ เพิ่ม
(3) ค่าฟังก์ชันวัตถุประสงค์จากตัวแก้ปัญหาตรงกับสูตร (\ref{eq:exp1-wass}) ที่ $\lambda \in \{0, 0.5, 1\}$
(4) ค่า LP คู่กันของวิธี B เท่ากับค่าที่ได้จากการหาค่าสูงสุดเหนือ $p$ โดยตรง ยืนยันข้อเสนอที่~\ref{prop:exp1-possB}
(5) $\kappa_E, \kappa_C$ ของวิธี A ตรงกับการจำลอง Monte Carlo
(6) น้ำหนักพอร์ตอยู่บน Simplex
(7) การทดสอบ Sharpe Ratio มีอัตราการปฏิเสธใกล้ $5\%$ เมื่อสมมติฐานหลักเป็นจริง

\subsection{ผลการทดลอง}

% ---- ตาราง/รูปด้านล่างสร้างโดย run_exp1.py ; ห้ามแก้ตัวเลขด้วยมือ ----
\begin{table}[htbp]
\centering
\caption{ผลนอกกลุ่มตัวอย่างที่ $\lambda = 0.5$ (ผลตอบแทนและ Sharpe Ratio รายปี, CVaR และ MDD เป็นร้อยละ)}
\label{tab:exp1-perf}
\resizebox{\linewidth}{!}{\input{results/real/table_perf_main.tex}}
\end{table}

\begin{table}[htbp]
\centering
\caption{คุณภาพของแนวหน้าพาเรโตนอกกลุ่มตัวอย่างและการทดสอบนัยสำคัญ}
\label{tab:exp1-front}
\resizebox{\linewidth}{!}{\input{results/real/table_front_main.tex}}
\end{table}

\begin{figure}[htbp]
\centering
\includegraphics[width=0.8\linewidth]{results/real/fig_oos_front.png}
\caption{แนวหน้าพาเรโตนอกกลุ่มตัวอย่างของแต่ละแบบจำลอง (หนึ่งจุดต่อหนึ่งค่า $\lambda$) เทียบกับพอร์ต $1/N$}
\label{fig:exp1-front}
\end{figure}

\begin{figure}[htbp]
\centering
\includegraphics[width=0.8\linewidth]{results/real/fig_wealth.png}
\caption{ความมั่งคั่งสะสมนอกกลุ่มตัวอย่างที่ $\lambda = 0.5$ แถบสีเทาคือช่วงตลาดผันผวนที่กำหนดล่วงหน้า}
\label{fig:exp1-wealth}
\end{figure}

\begin{table}[htbp]
\centering
\caption{CVaR$_{0.05}$ ของผลขาดทุนรายสัปดาห์ (ร้อยละ) แยกตามช่วงตลาด}
\label{tab:exp1-regime}
\resizebox{\linewidth}{!}{\input{results/real/table_regimes_cvar.tex}}
\end{table}

% ---- เขียนส่วนนี้หลังเห็นผลจริง ; ระบุทั้งผลที่สนับสนุนและไม่สนับสนุน ----
[ผลหลัก 1: เปรียบเทียบ Wasserstein กับ Nominal ด้าน CVaR, MDD และ HV
พร้อมช่วงความเชื่อมั่นจากตารางที่~\ref{tab:exp1-front}]

[ผลหลัก 2: เปรียบเทียบทุกแบบจำลองกับ $1/N$ หากค่า $p$ หลังปรับ Holm มากกว่า 0.05
ให้สรุปว่า ``ไม่พบหลักฐานว่าแตกต่าง'' ไม่ใช่ ``ดีกว่า'' หรือ ``แย่กว่า'']

[ผลหลัก 3: วิธี A ให้แนวหน้าสั้น ดูจากค่า \texttt{front\_spread.csv} ของ Poss-A เทียบกับ Poss-B
ซึ่งยืนยันข้อสังเกตในหัวข้อที่ 5.2 และแสดงว่าวิธี B แก้ปัญหานี้ได้หรือไม่]

[ผลหลัก 4: ช่วงตลาดผันผวน (ตารางที่~\ref{tab:exp1-regime}) แบบจำลองที่ทนทานลดผลขาดทุนส่วนหาง
ในช่วงใดบ้าง และเสียผลตอบแทนในช่วงปกติเท่าใด]

\subsection{ข้อจำกัดของการทดลองที่ 1}
\begin{enumerate}
\item วิธี B จำกัด Support ไว้ที่สถานการณ์จำลอง จึงเป็นค่าประมาณของ Possibility-induced
Ambiguity Set และบทแทรก 2.1 ที่รองรับวิธี A ยังอยู่ระหว่างการพิสูจน์ให้รัดกุม
\item เมทริกซ์ความผันผวนเป็นแบบแนวทแยงจาก GARCH รายสินทรัพย์ ความสัมพันธ์ระหว่างสินทรัพย์
มาจาก Residuals เชิงประจักษ์ซึ่งคงที่ภายในหน้าต่างข้อมูล
\item รัศมีที่ได้จากกฎ Bootstrap ไม่ได้ถูกปรับให้เหมาะกับผลการลงทุน ซึ่งเป็นความตั้งใจเพื่อป้องกัน
Data Snooping แต่ค่าที่เหมาะสมกว่าอาจมีอยู่ และจะศึกษาในการทดลองที่ 4
\item ผลเป็นของตลาดสหรัฐอเมริกาเท่านั้น การขยายไปยังตลาดไทยจะใช้ดัชนีหมวดอุตสาหกรรมของ
ตลาดหลักทรัพย์แห่งประเทศไทยด้วยขั้นตอนเดียวกัน
\end{enumerate}
